% R07 draft fragment; integrate only after checking the main manuscript's notation.
% Requires amsmath/amssymb/amsthm and theorem/proposition environments.
% Not a compiled or submitted manuscript.
\subsection{Kernel perturbations and cusp conditioning}
Let $F$ denote the polynomially deformed Jacobi cosine integral, and let
\[
 F_h(t;\lambda,\mu,\nu)=\int_0^\infty\Phi(u)(1+h(u))
 e^{\lambda u^2+\mu u^4+\nu u^6}\cos(2tu)\,du.
\]
Here $h$ is one fixed real measurable function, independent of all
controls, extended evenly to the real line. Put $D_n^h=\partial_t^nF_h$.

\begin{theorem}[Uniform local persistence]
For every $\|h\|_\infty\le10^{-18}$, the specified continuation tubes
over $-29\le\nu\le0$ contain one glued analytic cusp graph
$c_h(\nu)=(t_h,\lambda_h,\mu_h)$, with $D_3^h>0>D_4^h$ and
$\|c_h(\nu)-c(\nu)\|_\infty<4\cdot10^{-5}$.
Moreover $0.26<\mu_h'<0.34$, and the graph crosses $\mu=0$ exactly once.
In the coordinates $s=t-t_h$, $\ell=\lambda-\lambda_h$,
$m=\mu-\mu_h$, the local root classification holds throughout
\[
 |s|\le0.003,\qquad |\ell|\le10^{-6},\qquad |m|\le2\cdot10^{-9}.
\]
There are exactly three simple real zeros between the two folds and
one outside them off the discriminant; a double and a simple zero on
a fold; and one triple zero at the cusp. No zero crosses the $s$ endpoints.
For $0<\ell\le10^{-6}$, writing $W_h=m_{\rm high,h}-m_{\rm low,h}$,
\[
 \partial_\nu m_{\rm high,h}<0<\partial_\nu m_{\rm low,h},\qquad
 0.0003\ell^{3/2}<-\partial_\nu W_h<0.003\ell^{3/2}.
\]
These comparisons hold for each fixed $h$, in its own cusp-centered,
unrescaled control coordinates.
\end{theorem}

The proof uses $|D_n^h-D_n|\le10^{-18}B_n$, perturbed contraction
and residual bounds on 58 tubes, 57 root-containment joins, and new
cusp-centered Taylor and fold-transport bounds. A separate rational
implementation reconstructs the new Taylor enclosures, expanded
nonlinear polynomials, and fold-ODE estimates. Original integral
enclosures and positive moment majorants remain input assumptions.
The bound is sufficient and is not claimed optimal.

The cusp location can nevertheless be sensitive to selected directions.
For $h_\epsilon(u)=\epsilon\cos(2au)$, with $a$ the recorded exact
dyadic approximation to the quartic cusp coordinate $t_Q$, one has
\[
 F_{h_\epsilon}(t)=F(t)+\frac\epsilon2\{F(t+a)+F(t-a)\}.
\]
At $\epsilon=0$, fixed $\nu=0$, validated implicit differentiation gives
\[
 492590000000<\frac{d\mu_Q}{d\epsilon}<492610000000.
\]
This is an infinitesimal derivative, not a finite-displacement estimate.

For a complementary finite-amplitude construction, let
\[
 M_{nj}=\int_0^\infty\Phi(u)e^{\lambda_Qu^2+\mu_Qu^4}
 \cos(2ju)(2u)^n\cos(2t_Qu+n\pi/2)\,du,
 \quad j=1,2,3,4.
\]
Let $A=(M_{nj})_{n=0,1,2}$, $\alpha_j=(-1)^{j-1}\det A_{\widehat j}$,
$Z=\sum_j|\alpha_j|$, and $h_0(u)=\sum_j(\alpha_j/Z)\cos(2ju)$.
These definitions use the exact cusp and exact integrals.

\begin{proposition}[An exact pinned cusp]
The four cofactors have signs $(-,+,-,+)$, so $Z>0$ and
$\|h_0\|_\infty=1$. For every $|\epsilon|\le1/2$, the kernel
$\Phi(1+\epsilon h_0)$ is positive and even, and the exact point $Q$
remains a nondegenerate cusp in the $\nu=0$ section. Both $D_3$ and
$|D_4|$ there remain at least $0.7995$ times their original values.
\end{proposition}

Indeed $A\alpha=0$ by the alternating-minor identity, so the three
cusp equations vanish exactly for every amplitude. Forty rigorously
enclosed shifted-frequency integrals give
$|\sum_j(\alpha_j/Z)M_{nj}|<0.401|D_nF(Q)|$ for $n=3,4$.
This proves the claimed nondegeneracy and the nonzero $(\lambda,\mu)$
unfolding determinant. Rounded displayed coefficients do not have the
exact moment-annihilation property. This larger-amplitude proposition
does not assert persistence of the whole cusp arc, of the previous
explicit finite root window, or of its fold-motion direction.
